% Gemini theme
% See: https://rev.cs.uchicago.edu/k4rtik/gemini-uccs
% A fork of https://github.com/anishathalye/gemini

\documentclass[final]{beamer}

% ====================
% Packages
% ====================

\usepackage[T1]{fontenc}
\usepackage{lmodern}

% \usepackage[size=custom,width=120,height=72,scale=1.0]{beamerposter}
\usepackage[scale=1, size=a1, orientation=portrait]{beamerposter}
% \setlength{\paperwidth}{33.1in}
% \setlength{\paperheight}{23.4in}

\usetheme{gemini}
% \usecolortheme{uchicago}
\usecolortheme{stanford}
\usepackage{graphicx}
\usepackage{booktabs}
\usepackage{tikz}
\usepackage{pgfplots}
\pgfplotsset{compat=1.17}

% Added 
% \usepackage{cite}
\usepackage[sort&compress,numbers]{natbib}
\usepackage[backend=biber, style=ieee]{biblatex}
\addbibresource{bibliography.bib}
\renewcommand*{\bibfont}{\footnotesize}
% Circuitikz
\usepackage[american]{circuitikz}
\usepackage{placeins}
% \renewcommand{\section}{\FloatBarrier\Oldsection}

% Added by LF
\usepackage{bm}
% \usepackage{physics} % BraKet notation
% \usepackage{braket}
\usepackage{tikz}
\usepackage{pgfplots}
\usepgfplotslibrary{colorbrewer}
\pgfplotsset{compat=1.17}
\usetikzlibrary{fit}
\usetikzlibrary{arrows}
\usepackage{tikz-3dplot}
\usetikzlibrary{shapes.callouts}
\tikzset{
  level/.style   = { ultra thick, blue },
  connect/.style = { dashed, red },
  notice/.style  = { draw, rectangle callout, callout relative pointer={#1} },
  % cota/.style  = { latex-latex, black, thin, dashed },
  % cota/.style  = { latex-latex, very thin, black, dashed },
  % cota/.style  = { <->, very thin, black, dashed, latex-{latex[width=0.1mm]} },
  cota/.style  = {very thin, black, dashed, {Latex[length=0.5mm]}-{Latex[length=0.5mm]} },
  cotaa/.style  = {very thin, black, {Latex[length=0.5mm]}-{Latex[length=0.5mm]} },
  aw/.style  = { -latex, black, thin,  },
  fine/.style  = { very thin, dashed },
  finee/.style  = { thin, dashed, darkgray },
  labell/.style   = { text width=3cm, font=\sffamily\small}
}
\usepackage{empheq} % braces for subequations
\usepackage{cleveref} % cref{label}, Cref{label} to get the section that contains the exercise
\usepackage[american]{circuitikz} % american/european
% flat 3D plots
\usepgfplotslibrary{colormaps}
\pgfplotsset{compat=newest}
\pgfplotsset{
  colormap/twilight/.style={colormap={twilight}{[1pt]
  rgb(0pt)=(0.8857501584075443, 0.8500092494306783, 0.8879736506427196);
  rgb(25pt)=(0.38407269378943537, 0.46139018782416635, 0.7309466543290268);
  rgb(50pt)=(0.18488035509396164, 0.07942573027972388, 0.21307651648984993);
  rgb(75pt)=(0.6980608153581771, 0.3382897632604862, 0.3220747885521809);
  rgb(100pt)=(0.8857115512284565, 0.8500218611585632, 0.8857253899008712);
}}}

% This works both with inline lists and with macros containing lists
\newcommand*{\GetListMember}[2]{%
    \edef\dotheloop{%
    \noexpand\foreach \noexpand\a [count=\noexpand\i] in {#1} {%
        \noexpand\IfEq{\noexpand\i}{#2}{\noexpand\a\noexpand\breakforeach}{}%
    }}%
    \dotheloop
    \par%
}%

% https://tex.stackexchange.com/questions/184005/adding-arrows-to-a-gantt-chart
% GanttHeader setups some parameters for the rest of the diagram
% #1 Width of the diagram
% #2 Width of the space reserved for task numbers
% #3 Width of the space reserved for task names
% #4 Number of months in the diagram
% #5 Initial month number
% In addition to these parameters, the layout of the diagram is influenced
% by keys defined below, such as y, which changes the vertical scale
\def\GanttHeaderIntraYear#1#2#3#4#5{%
 \pgfmathparse{(#1-#2-#3)/#4}
 \tikzset{y=7mm, task number/.style={left, font=\bfseries},
     task description/.style={text width=#3,  right, draw=none,
           % font=\sffamily, xshift=#2,
           xshift=#2,
           minimum height=2em},
     gantt bar/.style={draw=black, fill=blue!30}, % default color
     help lines/.style={draw=black!50, dashed},
     x=\pgfmathresult pt
     }
  \def\totalmonths{#4}
  \def\monthsname{Jan,Feb,Mar,Apr,May,Jun,Jul,Aug,Sep,Oct,Nov,Dec}
  % \node (Header) [task description] at (0,0) {\textbf{\large Tasks}};
  \node (Header) [task description] at (0,0) {};
  \begin{scope}[shift=($(Header.south east)$)]
	  	\node[above, anchor=north] at (0,0.75) {\GetListMember{\monthsname}{1}};
	  \foreach[count = \z from #5+1, count = \x from #5] \y in {#5,...,#4}
	  	\node[above, anchor=north] at (\x,0.75) {\GetListMember{\monthsname}{\z}};
 \end{scope}
}


\def\GanttHeader#1#2#3#4#5{%
 \pgfmathparse{(#1-#2-#3)/#4}
 \tikzset{y=7mm, task number/.style={left, font=\bfseries},
     task description/.style={text width=#3,  right, draw=none,
           % font=\sffamily, xshift=#2,
           xshift=#2,
           minimum height=2em},
     gantt bar/.style={draw=black, fill=blue!30},
     help lines/.style={draw=black!50, dashed},
     x=\pgfmathresult pt
     }
  \def\totalmonths{#4}
  \def\monthsname{Jan,Feb,Mar,Apr,May,Jun,Jul,Aug,Sep,Oct,Nov,Dec}
  % \node (Header) [task description] at (0,0) {\textbf{\large Tasks}};
  \node (Header) [task description] at (0,0) {};
  \begin{scope}[shift=($(Header.south east)$)]
	  \foreach[count = \x from 0] \y in {#5,...,12}
	  	\node[above, anchor=north] at (\x,0.75) {\GetListMember{\monthsname}{\y}};
	  % \node[above] at (\x,1) {\Large{\y}};
    \foreach \x in {1,...,#4}
	  \node[above, anchor=north] at (\x,0.75) {\GetListMember{\monthsname}{\x}};
 \end{scope}
}


\newcommand\Task[5]{
	\Taskk{#1}{#2}{#3}{#4}{blue}{#5}
}


\def\Taskk#1#2#3#4#5#6{%
\node[task number] at ($(Header.west) + (0, -#1)$) {\normalfont#6};
\node[task description] at (0,-#1) {#2};
\begin{scope}[shift=($(Header.south east)$)]
  \draw (0,-#1) rectangle +(\totalmonths, 1);
  \foreach \x in {1,...,\totalmonths}
    \draw[help lines] (\x,-#1) -- +(0,1);
  \filldraw[gantt bar, fill=#5!30] ($(#3, -#1+0.2)$) rectangle +(#4,0.6);
\end{scope}
}

\def\Range#1#2#3#4#5#6{%
\node[task number] at ($(Header.west) + (0, -#1)$) {#6};
\node[task description] at (0,-#1) {\bfseries#2};
\begin{scope}[shift=($(Header.south east)$)]
  \draw (0,-#1) rectangle +(\totalmonths, 1);
  \foreach \x in {1,...,\totalmonths}
    \draw[help lines] (\x,-#1) -- +(0,1);
  \filldraw[gantt bar, fill=#5!30] ($(#3, -#1+0.4)$) rectangle +(#4,0.2);
\end{scope}
}





%%%%%%%%%%%%%%%%%%%%%%%
% Josephson junction

\makeatletter
\pgfcircdeclarebipolescaled{instruments}
{
    % put the node text above and centered
    \anchor{text}{\pgfextracty{\pgf@circ@res@up}{\northeast}
        \pgfpoint{-.5\wd\pgfnodeparttextbox}{
            \dimexpr.5\dp\pgfnodeparttextbox+.5\ht\pgfnodeparttextbox+\pgf@circ@res@up\relax
        }
    }
}
{\ctikzvalof{bipoles/oscope/height}}
{josephson}
{\ctikzvalof{bipoles/oscope/height}}
{\ctikzvalof{bipoles/oscope/width}}
{
    \pgf@circ@setlinewidth{bipoles}{\pgfstartlinewidth}
    \pgfextracty{\pgf@circ@res@up}{\northeast}
    \pgfextractx{\pgf@circ@res@right}{\northeast}
    \pgfextractx{\pgf@circ@res@left}{\southwest}
    \pgfextracty{\pgf@circ@res@down}{\southwest}
    \pgfmathsetlength{\pgf@circ@res@step}{0.25*\pgf@circ@res@up}
    \pgfscope
        \pgfpathrectanglecorners{\pgfpoint{\pgf@circ@res@left}{\pgf@circ@res@down}}{\pgfpoint{\pgf@circ@res@right}{\pgf@circ@res@up}}
        \pgf@circ@draworfill
    \endpgfscope
    \pgfscope
      \pgfpathmoveto{\pgfpoint{\pgf@circ@res@left}{\pgf@circ@res@up}}%
      \pgfpathlineto{\pgfpoint{\pgf@circ@res@right}{\pgf@circ@res@down}}%
      \pgfpathmoveto{\pgfpoint{\pgf@circ@res@right}{\pgf@circ@res@up}}%
      \pgfpathlineto{\pgfpoint{\pgf@circ@res@left}{\pgf@circ@res@down}}%
      \pgfusepath{draw}
    \endpgfscope
}
\def\pgf@circ@josephson@path#1{\pgf@circ@bipole@path{josephson}{#1}}
\tikzset{josephson/.style = {\circuitikzbasekey, /tikz/to path=\pgf@circ@josephson@path, l=#1}}



% Saturator block
\tikzset{%
  saturation block/.style={%
    draw,
    path picture={
      % Get the width and height of the path picture node
      \pgfpointdiff{\pgfpointanchor{path picture bounding box}{north east}}%
        {\pgfpointanchor{path picture bounding box}{south west}}
      \pgfgetlastxy\x\y
      % Scale the x and y vectors so that the range
      % -1 to 1 is slightly shorter than the size of the node
      \tikzset{x=\x*.4, y=\y*.4}
      %
      % Draw annotation
      \draw (-1,0) -- (1,0) (0,-1) -- (0,1);
      \draw (-1,-.7) -- (-.7,-.7) -- (.7,.7) -- (1,.7);
    }
  }
}



 % gantt diagram, Josephson junction and others
\usetikzlibrary{arrows}

\usepackage[caption=false]{subfig}

\usepackage[para]{threeparttable}

% \date{06th March 2023}
\date{\today}

% Added
\usepackage{siunitx}
\usepackage{cleveref}
\setlength{\tabcolsep}{0.3em} % for the horizontal table padding
\usepackage{pgfplots}
\usepgfplotslibrary{colorbrewer}
\usepackage{tikz}
\usetikzlibrary{plotmarks}
\usetikzlibrary{calc}
\usepackage[american]{circuitikz} % american/european
\usepackage{ifthen}
\usepackage{moresize}

% Reduce image separation with text
\setlength{\textfloatsep}{5pt}
% \usepackage[]{caption}
\captionsetup[figure]{skip=-10pt}
\setlength{\abovecaptionskip}{-4pt plus 1pt minus 1pt}

% Get citations in range, [3-6]
% \usepackage{cite}

% Pseudocode algorithms
\usepackage{algorithm}
\usepackage{algpseudocode}

\usepackage{anyfontsize}
\usepackage{svg}


% ====================
% Lengths
% ====================

% If you have N columns, choose \sepwidth and \colwidth such that
% (N+1)*\sepwidth + N*\colwidth = \paperwidth
\newlength{\sepwidth}
\newlength{\colwidth}
\setlength{\sepwidth}{0.025\paperwidth}
\setlength{\colwidth}{0.3\paperwidth}

\newcommand{\separatorcolumn}{\begin{column}{\sepwidth}\end{column}}


% \title{Your Title Goes Here}
\title{A Simple Zero Average Switching Energy Differential SAR ADC}

\author{Llorenç Fanals-i-Batllori \inst{1}\inst{,*} \and Namit Mishra \inst{1} \and Lorenzo Rota \inst{1} \and Aldo Pena-Perez \inst{1}}

\institute[shortinst]{\inst{1} SLAC National Accelerator Laboratory \samelineand \inst{*} \href{mailto:fanals@slac.stanford.edu}{fanals@slac.stanford.edu} }

% ====================
% Footer (optional)
% ====================

% \footercontent{
%    \href{rylanschaeffer.github.io}{rylanschaeffer.github.io} \hfill
%    \hfill
%    IEEE ISCAS 2024 - 19$^{\text{th}}$ to 22$^{\text{th}}$ May, Singapore \hfill
%   \href{mailto:rylanschaeffer@gmail.com}{rylanschaeffer@gmail.com}}
% % (can be left out to remove footer)

\footercontent{ \hfill IEEE ISCAS 2024 - May 19 to 22, Singapore \hfill}

% ====================
% Logo (optional)
% ====================

% use this to include logos on the left and/or right side of the header:
% \logoright{\includegraphics[height=7cm]{logos/cs-logo-maroon.png}}
% \logoleft{\includegraphics[height=7cm]{logos/cs-logo-maroon.png}}
\logoright{\includegraphics[height=6cm]{stanford_logos/iscas24b.png}}
\logoleft{\includegraphics[height=4cm]{stanford_logos/logo_SLAC.png}}

% ====================
% Body
% ====================

\begin{document}

% % This adds the Logos on the top left and top right
% \addtobeamertemplate{headline}{}
% {
%     \begin{tikzpicture}[remember picture,overlay]
%     %   \node [anchor=north west, inner sep=3cm] at ([xshift=0.0cm,yshift=1.0cm]current page.north west)
%     %   {\includegraphics[height=5.0cm]{stanford_logos/Stanford-CS.png}}; % uc-logo-white.eps
%       \node [anchor=north east, inner sep=3cm] at ([xshift=0.0cm,yshift=2.5cm]current page.north east)
%       {};
%     \end{tikzpicture}
% }

\begin{frame}[t]
\begin{columns}[t]
%%%%%%%%%%%%%%%%%%%%%%%%%%%%%%%%%%%%%%%%%%%%%%%%%%%%%%%%%%%%%%%%%%%%%%%%%%%%%%%%%%%%%%%%%%
\separatorcolumn
%%%%%%%%%%%%%%%%%%%%%%%%%%%%%%%%%%%%%%%%%%%%%%%%%%%%%%%%%%%%%%%%%%%%%%%%%%%%%%%%%%%%%%%%%%

\begin{column}{\colwidth}

  \begin{block}{Introduction}
      % The successive approximation register (SAR) analog-to-digital converter (ADC) has emerged as a dominant architecture for low-power applications. This is attributed to its superior energy efficiency, compact form factor, and seamless compatibility with technology scaling. Nevertheless, recent advancements have notably focused on refining the switching logic to surpass the existing limits in power consumption. Past works such as \cite{Yuan12, Sanyal13, Zhu13, Sanyal13b, Chen14} have demonstrated energy saving beyond the \SI{95}{\percent} with respect to the conventional switching but at the expense of variable common-mode voltages at the comparator inputs. As a consequence, a special comparator design is necessary, restricting the effective resolution. Alternative solutions have addressed this limitation; however, they come with trade-offs such as doubling the capacitor area \cite{Zhu10}, incurring substantial switching energy \cite{Rahimi14}, or halving the input signal amplitude \cite{Xin21}. The methods in \cite{Ding18} and \cite{Yu23} leverage bridge capacitor designs to achieve significant area reduction, yet implementing this solution requires careful layout to minimize undesired couplings around the attenuation capacitor which could otherwise limit the linearity. Remarkable energy savings exceeding \SI{98}{\percent} have been demonstrated in prior works \cite{Wu17, Li20, Hu22, Momeni20} while maintaining common-mode voltage variations and overall performance; however, with recent advancements, it is now achievable to reach \SI{100}{\percent} energy savings. An example of this progress can be found in \cite{Akbari19}; although the capacitive DAC employed in this approach is larger compared to the other techniques discussed here. Similarly, \cite{Yousefi18} accomplishes zero average energy and a reduced capacitance area of \SI{87}{\percent}; nonetheless, the implementation of the switching network and the associated digital logic is complex and requires extra clock cycles to perform the conversion.

% The switching scheme presented here achieves a zero average switching energy with a fully-differential input signal and requires only one clock period per resolved bit. While architectures \cite{Zhu10,Rahimi14,Xin21} and \cite{Wu17,Li20,Hu22}  claim to maintain a constant common-mode voltage, they actually experience a variation when discriminating the least significant bit (LSB). In contrast, the proposed solution introduces a marginal shift for only two specific digital codes during the last transition, representing less than \SI{1}{\percent} of the total amount of cases when considering reasonable resolutions (i.e.,\,$\geq$\,9\,bit).

The successive approximation register (SAR) analog-to-digital converter (ADC) has emerged as a dominant architecture for low-power applications. This is attributed to its superior energy efficiency, compact form factor, and seamless compatibility with technology scaling. Recent advancements have notably focused on refining the switching logic. Remarkable energy savings exceeding \SI{98}{\percent} have been demonstrated in prior works \cite{Wu17, Li20, Hu22, Momeni20} while maintaining common-mode voltage variations and overall performance. However, with recent advancements, it is now achievable to reach \SI{100}{\percent} energy savings at the cost of a large capacitive DAC as in \cite{Akbari19} or complexity and extra clock cycles as in \cite{Yousefi18}.    
% 
The proposed switching scheme achieves a zero average switching energy with a fully-differential input signal, an almost constant common-mode voltage and requires only one clock period per resolved bit.

  \end{block}


  \begin{block}{Proposed switching scheme}

    \begin{itemize}
      \item Symmetric operation.
      \item Zero energy during successive approximations.
      \item Almost constant common-mode voltage.
      \item Compact capacitive DAC.
      \item No additional clock cycles.
    \end{itemize}
    \vspace{-1cm}

      \begin{figure} \centering
              \centering
      	    \newcommand{\state}[8]{
	% #1, #2: x and y offset position
	% #3, #4, #5: top voltage values, from left to right
	% #6, #7, #8: bot voltage values, from left to right
	\draw ($ ({#1},{#2}) + (0.65,-2.55)$) node[] { $C$};
	\draw ($ ({#1},{#2}) + (0.15,-2.55)$) node[] { $C$};
	\draw ($ ({#1},{#2}) + (-0.3,-2.55)$) node[] { $2C$};
	\draw ($ ({#1},{#2}) + (-0.5,-2.15)$) to[C] ($ ({#1},{#2}) + (-0.5,-2.65)$);
	\draw ($ ({#1},{#2}) + (0.0,-2.15) $) to[C] ($ ({#1},{#2}) + (0.0,-2.65)$);
	\draw ($ ({#1},{#2}) + (0.5,-2.15) $) to[C] ($ ({#1},{#2}) + (0.5,-2.65)$);
	\draw ($ ({#1},{#2}) + (-0.5,-2.65)$) to[short] ($ ({#1},{#2}) + (0.5,-2.65)$);

	\draw ($ ({#1},{#2}) + (1.25,-2.825)$) node[plain amp, line width=0.2, yscale=-1] (opamp1) {};
	\draw ($ ({#1},{#2}) + (1.08,-3.0) $) node[] { $-$};
	\draw ($ ({#1},{#2}) + (1.08,-2.65) $)node[] { $+$};
	\draw (opamp1.-) to[short] ($ ({#1},{#2})+(0.5,-3.0) $);
	\draw (opamp1.+) to[short] ($ ({#1},{#2})+(0.5,-2.65)$);
	
	\draw ($ ({#1},{#2}) + (0.65,-3.1)$) node[] { $C$};
	\draw ($ ({#1},{#2}) + (0.15,-3.1)$) node[] { $C$};
	\draw ($ ({#1},{#2}) + (-0.3,-3.1)$) node[] { $2C$};
	\draw ($ ({#1},{#2}) + (-0.5,-3.0)$) to[C] ($ (#1,#2)+(-0.5,-3.5)$);
	\draw ($ ({#1},{#2}) + (0.0,-3.0) $) to[C] ($ (#1,#2)+(0.0,-3.5)$);
	\draw ($ ({#1},{#2}) + (0.5,-3.0) $) to[C] ($ (#1,#2)+(0.5,-3.5)$);
	\draw ($ ({#1},{#2}) + (-0.5,-3.0)$) to[short] ($ (#1,#2)+(0.5,-3.0)$);


	\ifthenelse{\equal{#3}{C}}{
		\draw[gray] ($ ({#1},{#2}) + (-0.5,-3.5) $) to[short] ($ ({#1},{#2}) + (-0.75,-3.5) $);
		\draw[gray] ($ ({#1},{#2}) + (-0.75,-3.5)$) to[short] ($ ({#1},{#2}) + (-0.75,-2.15)$);
		\draw[gray] ($ ({#1},{#2}) + (-0.5,-2.15)$) to[short] ($ ({#1},{#2}) + (-0.75,-2.15)$);
	}{
		\draw ($ ({#1},{#2}) + (-0.5,-2.0)$) node[]  { $#3$};
		\draw ($ ({#1},{#2}) + (-0.5,-3.65)$) node[] { $#6$};
	}

	\ifthenelse{\equal{#4}{C}}{
		\draw[gray] ($ ({#1},{#2}) + (-0.0,-3.5) $) to[short] ($ ({#1},{#2}) + (-0.25,-3.5) $);
		\draw[gray] ($ ({#1},{#2}) + (-0.25,-3.5)$) to[short] ($ ({#1},{#2}) + (-0.25,-2.15)$);
		\draw[gray] ($ ({#1},{#2}) + (-0.0,-2.15)$) to[short] ($ ({#1},{#2}) + (-0.25,-2.15)$);
	}{
		\draw ($ ({#1},{#2}) + (0.0,-2.0) $) node[]  { $#4$};
		\draw ($ ({#1},{#2}) + (0.0,-3.65) $) node[] { $#7$};
	}

	\ifthenelse{\equal{#5}{C}}{
		\draw[gray] ($ ({#1},{#2}) + (0.5,-3.5)  $) to[short] ($ ({#1},{#2}) + (0.25,-3.5)  $);
		\draw[gray] ($ ({#1},{#2}) + (0.25,-3.5) $) to[short] ($ ({#1},{#2}) + (0.25,-2.15) $);
		\draw[gray] ($ ({#1},{#2}) + (0.5,-2.15) $) to[short] ($ ({#1},{#2}) + (0.25,-2.15) $);
	}{
		\draw ($ ({#1},{#2}) + (0.5,-2.0) $) node[]  { $#5$};
		\draw ($ ({#1},{#2}) + (0.5,-3.65) $) node[] { $#8$};
	}

	% \ifthenelse{\equal{#3}{C}}{
	% 	\draw (1,1) node[] {A};
	% }{
	% 	\draw (1,1) node[] {B};
	% }
}

\newcommand{\labels}[4]{
	% #1, #2: x and y offset position
	% #3, #4: top and bot voltage increments, respectively

	% \draw ($ ({#1},{#2}) + (0.5,-2.75) $) node[] {\tiny $\Delta = {#3}$};
	% \draw ($ ({#1},{#2}) + (-0.25,-2.9) $) node[] {\tiny $\Delta = {#4}$};
	\ifthenelse{\equal{#3}{}}{
	}{
		\draw ($ ({#1},{#2}) + (0.25,-2.75) $) node[] { $\Delta = {#3}$};
	}
	\ifthenelse{\equal{#4}{}}{
	}{
		\draw ($ ({#1},{#2}) + (-0.25,-2.9) $) node[] { $\Delta = {#4}$};
	}

}

\newcommand{\switches}[2]{
	% #1, #2: x and y offset position
	\draw ($ ({#1},{#2}) + (-1.0,-2.65) $) to[ncs] ($ ({#1},{#2}) + (-0.5,-2.65) $);
	\draw ($ ({#1},{#2}) + (-1.0,-3.00) $) to[ncs] ($ ({#1},{#2}) + (-0.5,-3.00) $);

	\draw ($ ({#1},{#2}) + (-1.2,-2.50) $) node[] { $v_{i,top}$};
	\draw ($ ({#1},{#2}) + (-1.2,-3.10) $) node[] { $v_{i,bot}$};
}


% Schematic created with CircuiTikZ-Tool, www.github.com/LFanals/CircuiTikZ-Tool.git 
\begin{circuitikz}[american, /tikz/circuitikz/bipoles/length=0.5cm, scale=1.3, transform shape, /tikz/circuitikz/bipoles/thickness=1]
\tikzstyle{block} = [draw, rectangle, minimum height=1cm, minimum width=2cm]
% \tikzstyle{every node}=[font=\tiny]
\tikzstyle{every node}=[font=\fontsize{7}{7}\selectfont]

\switches{5.5}{2}
\state{5.5}{2}{1}{0}{0}{1}{0}{0}        \labels{5.5}{2}{}{}

\draw (5.8,-0.8) node[] { $v_{top}$};
\draw (5.3,-0.9) node[] { $v_{bot}$};
 
% top side
\state{9}{6}{0.5}{0}{0}{1}{0.5}{0.5}  \labels{9}{6}{- 1/4}{+ 1/4}

\state{12.5}{8}{0}{0.5}{0}{1}{1}{0.5} \labels{12.5}{8}{- 1/8}{+ 1/8}
\state{12.5}{4}{C}{0}{0}{C}{0.5}{0.5} \labels{12.5}{4}{+ 1/8}{- 1/8}

\state{16}{9}{0}{0.5}{0}{1}{1}{1}     \labels{16}{9}{0}{+ 1/8}
\state{16}{7}{0}{C}{0}{1}{C}{0.5}     \labels{16}{7}{+ 1/16}{- 1/16}
\state{16}{5}{0.5}{C}{0}{1}{C}{0.5}   \labels{16}{5}{- 1/16}{+ 1/16}
\state{16}{3}{C}{C}{0}{C}{C}{0.5}     \labels{16}{3}{+ 1/16}{- 1/16}

% bot side
\state{9}{-2}{1}{0.5}{0.5}{0.5}{0}{0}  \labels{9}{-2}{+ 1/4}{- 1/4}

\state{12.5}{-4}{1}{1}{0.5}{0}{0.5}{0} \labels{12.5}{-4}{+ 1/8}{- 1/8}
\state{12.5}{0}{C}{0.5}{0.5}{C}{0}{0} \labels{12.5}{0}{- 1/8}{+ 1/8}

\state{16}{-5}{1}{1}{1}{0}{0.5}{0}     \labels{16}{-5}{+ 1/8}{0}
\state{16}{-3}{1}{C}{0.5}{0}{C}{0}     \labels{16}{-3}{- 1/16}{+ 1/16}
\state{16}{-1}{1}{C}{0.5}{0.5}{C}{0}   \labels{16}{-1}{+ 1/16}{- 1/16}
\state{16}{1}{C}{C}{0.5}{C}{C}{0}     \labels{16}{1}{- 1/16}{+ 1/16}


\begin{scope}[xshift=0cm, yshift=-0.325cm]
	\draw [->] (14.5,5.5) -- node[midway,left] {$E=0$} (15.0,6.5);
	\draw [->] (14.5,5.5) -- node[midway,left] {$E=5/32$} (15.0,4.5);
	\draw [->] (14.5,1.5) -- node[midway,left] {$E=5/32$} (15.0,2.5);
	\draw [->] (14.5,1.5) -- node[midway,left] {$E=1/32$} (15.0,0.5);
	\draw [->] (11.0,3.5) -- node[midway,left] {$E=3/8$} (11.5,5.5);
	\draw [->] (11.0,3.5) -- node[midway,left] {$E=1/8$} (11.5,1.5);
	\draw [->] (7.5,-0.5) -- node[midway,left] {$E=-1/2$} (8,3.5);
\end{scope}

\begin{scope}[xshift=0cm, yshift=-8.325cm]
	\draw [->] (14.5,5.5) -- node[midway,left] {$E=1/32$} (15.0,6.5);
	\draw [->] (14.5,5.5) -- node[midway,left] {$E=5/32$} (15.0,4.5);
	\draw [->] (14.5,1.5) -- node[midway,left] {$E=5/32$} (15.0,2.5);
	\draw [->] (14.5,1.5) -- node[midway,left] {$E=0$} (15.0,0.5);
	\draw [->] (11.0,3.5) -- node[midway,left] {$E=1/8$} (11.5,5.5);
	\draw [->] (11.0,3.5) -- node[midway,left] {$E=3/8$} (11.5,1.5);
	\draw [->] (7.5,7.5) -- node[midway,left]  {$E=-1/2$} (8,3.5);
\end{scope}


\end{circuitikz}

   

          	\caption{Proposed symmetric switching scheme for a 4-bit SAR ADC. Top sampling is performed. Voltage increments following each transition are visualized at both the top and bottom nodes. The energy values are indicated on the transition arrows, relative to $C V^2 = C$. Up transitions occur when $v_{top} > v_{bot}$.}
              \label{fig:scheme}
      \end{figure}

  \end{block}

  \begin{block}{Proposed switching scheme algorithm}

      \begin{figure} \centering
              \centering
      	    
\begin{circuitikz}[american, /tikz/circuitikz/bipoles/length=0.5cm, scale=1.5, transform shape, /tikz/circuitikz/bipoles/thickness=1]
    \tikzstyle{every node}=[font=\fontsize{6}{6}\selectfont]
% Algorithm
\begin{scope}[xshift=0.2cm]
	% \draw [->] (18,4) -- node[midway,left] {$i == N-2$} (18,3.5);
	\draw [->] (19,4.65) -- node[midway,left] {} (19,3.27);
	\draw [->] (19,1.87) -- node[midway,left] {} (19,0.25);
    \draw[->, bend right=90, distance=0.75cm] (19,-6.8) to node[midway, left, yshift=-0.2cm, xshift=0.6cm] { if $i < N-2$\ :} (20,-6.8);
    \draw (19.4,-7.9) node[] { $i \leftarrow i + 1$};
	\draw [-] (21,-6.8) |- node[] {} (23.25,-7.25); % COMPARE to INITIALIZE
    \draw [-] (23.25,-7.25) -- node[xshift=-0.75cm, yshift=3cm] { if $i==N-2$} (23.25,5.0); % COMPARE to INITIALIZE
	\draw [->] (23.25,5.0) |- node[midway,left] {} (22.05,5.5); % COMPARE to INITIALIZE

% INITIALIZE
\begin{scope}[xshift=18cm, yshift=6cm]
    \path [fill=gray!10, draw=black, line width=0.5pt] (0,0.25) rectangle (4,-1.35);
    \draw (0, 0) node[right] { \textbf{INITIALIZATION}};
    \draw (0, -0.3) node[right] { Sampling switches $\leftarrow$ open};
    \draw (0, -0.6) node[right] { $V_{a} \leftarrow [1,0,...,0], V_{b} \leftarrow [1,0,...,0]$};
    \draw (0, -0.9) node[right] { $V_{top} \leftarrow \text{flag}_\text{top} \ ? \ V_a : V_b$};
    \draw (0, -1.2) node[right] { $V_{bot} \leftarrow \text{flag}_\text{top} \ ? \ V_b : V_a$};
\end{scope}
% SAMPLE
\begin{scope}[xshift=18cm, yshift=3cm]
    \path [fill=gray!10, draw=black, line width=0.5pt] (0,0.25) rectangle (3.0,-1.1);
    \draw (0, 0) node[right] { \textbf{SAMPLING}};
    \draw (0, -0.3) node[right] { Sampling switches $\leftarrow$ close};
    \draw (0, -0.6) node[right] { $D \leftarrow 0$, $i \leftarrow 0$};
    \draw (0, -0.9) node[right] { $\text{flag}_{\text{extreme}} \leftarrow 1$};
\end{scope}
% COMPARE
\begin{scope}[xshift=18cm, yshift=0cm]
    \path [fill=gray!10, draw=black,line width=0.5pt] (0,0.25) rectangle (5.0,-6.8);
    \draw (0, 0) node[right] { \textbf{SUCCESSIVE APPROXIMATIONS}};
    \draw (0, -0.3) node[right] { $D \leftarrow D + (v_{top} > v_{bot}) << (N - 1 - i)$};
    % i = 0
    \draw (0, -0.6) node[right] { if $i == 0$:};
    \draw (0.3, -0.9) node[right] { Sampling switches $\leftarrow$ open};
    \draw (0.3, -1.2) node[right] { $(v_{top} > v_{bot}) \ ? \ \text{flag}_\text{top} \leftarrow 1 : \text{flag}_\text{top} \leftarrow 0 $};
    \draw (0.3, -1.5) node[right] { $V_a[0] \leftarrow 0.5$};
    \draw (0.3, -1.8) node[right] { $V_b[1:N-2] \leftarrow [0.5,0.5,...,0.5]$};
    % i<= N-2
    \draw (0, -2.1) node[right] { else if $i <= N-2$:};
        % vtop > vbot
    \draw (0.3, -2.4) node[right] { if $v_{top} > v_{bot}$:};
            % flag_extreme
    \draw (0.6, -2.7) node[right] { if $\text{flag}_\text{extreme}$:};
    \draw (0.9, -3.0) node[right] { if $i < N - 2: \ V_a[i-1,i] \leftarrow [0,0.5]$};
    \draw (0.9, -3.3) node[right] { $V_b[i] \leftarrow 1$};
            % !flag_extreme
    \draw (0.6, -3.6) node[right] { else:};
    \draw (0.9, -3.9) node[right] { $V_{a,prev}[i-1] \leftarrow V_a[i-1]$}; % store
    \draw (0.9, -4.2) node[right] { $V_{b,prev}[i-1] \leftarrow V_b[i-1]$}; % store
    \draw (0.9, -4.5) node[right] { $V_a[i-1] \leftarrow M, V_b[i-1] \leftarrow M$}; % merge
    \draw (0.9, -4.8) node[right] { $V_a[i-2] \leftarrow V_{a,prev}[i-2]$}; % load
    \draw (0.9, -5.1) node[right] { $V_b[i-2] \leftarrow V_{b,prev}[i-2]$}; % load
        % else
    \draw (0.3, -5.4) node[right] { else:};
    \draw (0.6, -5.7) node[right] { $V_{a,prev}[i-1] \leftarrow V_a[i-1]$}; % store
    \draw (0.6, -6.0) node[right] { $V_{b,prev}[i-1] \leftarrow V_b[i-1]$}; % store
    \draw (0.6, -6.3) node[right] { $V_a[i-1] \leftarrow M, V_b[i-1] \leftarrow M$}; % merge
    % if (vtop > vbot XOR flag_top): flag_extreme <- 0
    \draw (0.3, -6.6) node[right] { if (($v_{top} > v_{bot}) \oplus \text{flag}_\text{top}$) \ : $\text{flag}_\text{extreme} \leftarrow 0$};
\end{scope}
\end{scope} % end algorithm
\end{circuitikz}



              \caption{Proposed switching scheme algorithm in pseudocode. $D$ is the resolved code and $i$ the comparison counter that ranges from $0$ to $N-1$. The $<<$ symbol symbolizes the shift left operation. $M$ indicates merging with its counterpart capacitor. Each arrow symbolizes a transition that occurs at the clock edge.}
              \label{fig:algo}
      \end{figure}

  \end{block}

  

\end{column}

%%%%%%%%%%%%%%%%%%%%%%%%%%%%%%%%%%%%%%%%%%%%%%%%%%%%%%%%%%%%%%%%%%%%%%%%%%%%%%%%%%%%%%%%%%
\separatorcolumn
%%%%%%%%%%%%%%%%%%%%%%%%%%%%%%%%%%%%%%%%%%%%%%%%%%%%%%%%%%%%%%%%%%%%%%%%%%%%%%%%%%%%%%%%%%

\begin{column}{\colwidth}

\begin{block}{Switching energy general observations}

    The energy consumed during the switching process in one capacitor array half $k$ (either top or bottom) between two successive steps, as can be generalized from the energy calculations in \cite{Ginsburg}, is as follows:
    \begin{equation} \label{eq:e}
        E_{k} = C V_{0,k} \left( \Delta V_{0,k} - \Delta v_{k} \right) + \sum_{i=0}^{N-3} 2^i C V_{i+1,k} \left( \Delta V_{i+1,k} - \Delta v_{k} \right) \ .
    \end{equation}
    Some main design strategies can be derived from this: 
    \begin{enumerate}
      \item Operating multiple switches simultaneously may actually result in lower energy consumption compared to using a single switch \cite{Xin21}.
      \item Removing the largest capacitors on the  CDAC has a notable impact on energy usage, particularly during the initial conversions due to the binary weighting configuration\,\cite{Yuan12, Sanyal13, Zhu13}. 
      \item If $V_{i,\text{k}} = 0$ (i.e., ground or negative reference), its contribution to the total switching energy becomes negligible. 
      \item When capacitor terminals are merged in a specific sequence as demonstrated in \cite{Xin21} and \cite{Wu17}, their connection becomes independent of voltage references, resulting in zero energy.
    \end{enumerate}

    \begin{figure}[t] \centering
    	\begin{circuitikz}[american, scale=1, transform shape]
    % \tikzstyle{block} = [draw, rectangle, minimum height=1cm, minimum width=2cm]
    % \tikzstyle{every node}=[font=\large]
    \tikzstyle{every node}=[font=\footnotesize]

    \draw ($ (-9.5,0.5) + (0,0) $) to[nos] ($ (-7.5,0.5) + (0,0) $);
    \draw (-9.5,0.75) node[] {$v_{i,top}$};
    \draw ($ (-9.5,-0.5) + (0,0) $) to[nos] ($ (-7.5,-0.5) + (0,0) $);
    \draw (-9.5,-0.75) node[] {$v_{i,bot}$};

    
\draw (1.5,-0.0) node[op amp,scale=1.02,yscale=-1] (opamp1) {};
\draw (opamp1.-) to[short] (-1.0,-0.5);
\draw (opamp1.+) to[short] (-1.0,0.5);

\draw (-0.5,0.5) to[C,l_=$C$] (-0.5,2.5);
\draw (-2.5,0.5) to[C,l_=$C$] (-2.5,2.5);
\draw (-4.5,0.5) to[C,l_=$2C$] (-4.5,2.5);
\draw (-7.5,0.5) to[C,l_=$2^{N-3}C$] (-7.5,2.5);
\draw (-0.5,2.8) node[] {$V_{0,top}$};
\draw (-2.5,2.8) node[] {$V_{1,top}$};
\draw (-4.5,2.8) node[] {$V_{2,top}$};
\draw (-7.5,2.8) node[] {$V_{N-2,top}$};
\draw (-7.5,0.5) to[short] (-0.5,0.5);
\node at (-7.5,2.5) [circle, fill, inner sep=1.5pt] {};
\node at (-4.5,2.5) [circle, fill, inner sep=1.5pt] {};
\node at (-2.5,2.5) [circle, fill, inner sep=1.5pt] {};
\node at (-0.5,2.5) [circle, fill, inner sep=1.5pt] {};
\draw (-5.4,1.5) node[] {$\cdots$};
\draw (-5.5,0.25) node[] {$v_{top}$};

\draw (-0.5,-0.5) to[C,l=$C$] (-0.5,-2.5);
\draw (-2.5,-0.5) to[C,l=$C$] (-2.5,-2.5);
\draw (-4.5,-0.5) to[C,l=$2C$] (-4.5,-2.5);
\draw (-7.5,-0.5) to[C,l=$2^{N-3}C$] (-7.5,-2.5);
\draw (-0.5,-2.8) node[] {$V_{0,bot}$};
\draw (-2.5,-2.8) node[] {$V_{1,bot}$};
\draw (-4.5,-2.8) node[] {$V_{2,bot}$};
\draw (-7.5,-2.8) node[] {$V_{N-2,bot}$};
\draw (-7.5,-0.5) to[short] (-0.5,-0.5);
\node at (-0.5,-2.5) [circle, fill, inner sep=1.5pt] {};
\node at (-2.5,-2.5) [circle, fill, inner sep=1.5pt] {};
\node at (-4.5,-2.5) [circle, fill, inner sep=1.5pt] {};
\node at (-7.5,-2.5) [circle, fill, inner sep=1.5pt] {};
\draw (-5.4,-1.5) node[] {$\cdots$};
\draw (-6.0,-0.25) node[] {$v_{bot}$};

\end{circuitikz}



% \begin{circuitikz}[ american, scale=0.6, transform shape]
% \tikzstyle{block} = [draw, rectangle, minimum height=1cm, minimum width=2cm]
% \draw (1.5,-9.0) node[op amp,scale=1.02,yscale=-1] (opamp1) {};
% \draw (opamp1.+) to[short] (0.0,-8.5);
% \draw (opamp1.-) to[short] (0.0,-9.5);
% \draw (0.0,-9.5) to[C] (0.0,-11.0);
% \draw (-1.0,-9.5) to[C] (-1.0,-11.0);
% \draw (-2.0,-9.5) to[C] (-2.0,-11.0);
% \draw (0.0,-8.5) to[C] (0.0,-7.0);
% \draw (-1.0,-8.5) to[C] (-1.0,-7.0);
% \draw (-2.0,-8.5) to[C] (-2.0,-7.0);
% \draw (0.0,-8.5) to[short] (-2.0,-8.5);
% \draw (0.0,-9.5) to[short] (-2.0,-9.5);
% \draw (-2.0,-6.6) node[] {$V_{0,\text{top}}$};
% \draw (-1.0,-6.6) node[] {$V_{1,\text{top}}$};
% \draw (0.0,-6.6) node[] {$V_{2,\text{top}}$};
% \draw (-2.0,-11.4) node[] {$V_{0,\text{bot}}$};
% \draw (-1.0,-11.4) node[] {$V_{1,\text{bot}}$};
% \draw (0.0,-11.4) node[] {$V_{2,\text{bot}}$};
% 
% \node at (-2.0,-7.0) [circle, fill, inner sep=1.5pt] {};
% \node at (-1.0,-7.0) [circle, fill, inner sep=1.5pt] {};
% \node at (-0.0,-7.0) [circle, fill, inner sep=1.5pt] {};
% 
% \node at (-2.0,-11.0) [circle, fill, inner sep=1.5pt] {};
% \node at (-1.0,-11.0) [circle, fill, inner sep=1.5pt] {};
% \node at (-0.0,-11.0) [circle, fill, inner sep=1.5pt] {};
% 
% \draw (-1.8,-8.8) node[] {$v_{\text{top}}$};
% \draw (-1.8,-9.3) node[] {$v_{\text{bot}}$};
% \begin{scope}[xshift=-7cm]
%     \draw (7.3,-9.85) node[] {$C$};
%     \draw (6.3,-9.85) node[] {$C$};
%     \draw (5.3,-9.85) node[] {$2C$};
%     \draw (7.3,-8.25) node[] {$C$};
%     \draw (6.3,-8.25) node[] {$C$};
%     \draw (5.3,-8.25) node[] {$2C$};
% \end{scope}
% \begin{scope}[yshift=0.85cm, xshift=-7cm]
%     \draw (7.5,-8.15) node[] {\small$C_{2,\text{top}}$};
%     \draw (6.5,-8.15) node[] {\small$C_{1,\text{top}}$};
%     \draw (5.5,-8.15) node[] {\small$C_{0,\text{top}}$};
% \end{scope}
% \begin{scope}[yshift=-2.6cm, xshift=-7cm]
%     \draw (7.5,-8.15) node[] {\small$C_{2,\text{bot}}$};
%     \draw (6.5,-8.15) node[] {\small$C_{1,\text{bot}}$};
%     \draw (5.5,-8.15) node[] {\small$C_{0,\text{bot}}$};
% \end{scope}
% \end{circuitikz}


% \begin{circuitikz}[ american, scale=0.6, transform shape]
% \tikzstyle{block} = [draw, rectangle, minimum height=1cm, minimum width=2cm]
% \draw (1.5,-9.0) node[op amp,scale=1.02,yscale=-1] (opamp1) {};


    	\caption{Generic $N$-bit SAR ADC with labeled binary weighted capacitors and voltages that are set by the SAR logic.}
        \label{fig:labels}
    \end{figure}


\end{block}


\begin{block}{Proposed scheme switching energy}
    For the specific cases of $N=4$ and $N=5$, the transition energies are demonstrated through the following summations:
    \begin{equation} \label{eq:en}
    	\begin{aligned}
    		& E(N=4) = C V^2 \left( - \frac{1}{2} + \frac{1}{4} + \frac{11}{128} \right), \\
    		& E(N=5) = C V^2 \left( -1 + \frac{1}{2} + \frac{1}{4} + \frac{55}{512} \right). \\
    	\end{aligned}
    \end{equation}
    As can be noticed from \eqref{eq:en}, the presence of a negative term in the first transitions plays a crucial role in minimizing the switching energy. By re-writing the equations in \eqref{eq:en} as a function of $N$, the overall switching energy for an $N$-bit resolution is found to be:
    \begin{equation} \label{eq:et}
    	E(N) = - C V^2 \left( 2^{-3} + 2^{-N-1} + 2^{-2N+1} \right) . 
    \end{equation}
    Negative energies $E(N)$ are possible when there is a positive switching energy.

\end{block}


\begin{block}{Behavioral simulations}
    \begin{figure} \centering
       \begin{tikzpicture}
    	   % red, blue, green, darkgray, orange, purple, brown, black, gray, olive, teal
        \begin{axis}[width=16cm, height=10cm, legend style={at={(0,1.0)}, legend columns=5, anchor=south west}, ylabel={\textit{Amplitude (\SI{}{\V})}}, xlabel={\textit{Time (\SI{}{\us})}}, grid=minor, enlarge x limits=false, every axis y label/.style={at={(ticklabel cs:0.5)},rotate=90,anchor=near ticklabel}, x filter/.code={\pgfmathmultiply{#1}{1000000}}]
    
            \addplot[mark=none, mark size = 1.5pt, line width=1.0pt, color=black] table [each nth point=1, x={t}, y={vcip_pseudo}, col sep=comma] {data_simulator_new/ADC_S3_waveforms_veriloga.csv};
            \addlegendentry{\scriptsize{CDAC} $v_{top}$}
            \addplot[mark=none, mark size = 1.5pt, line width=1.0pt, color=gray] table [each nth point=1, x={t}, y={vcin_pseudo}, col sep=comma] {data_simulator_new/ADC_S3_waveforms_veriloga.csv};
            \addlegendentry{\scriptsize{CDAC} $v_{bot}$}
            \addplot[mark=none, dashed, mark size = 1.5pt, line width=0.75pt, color=black] table [each nth point=1, x={t}, y={vinp}, col sep=comma] {data_simulator_new/ADC_S3_waveforms_veriloga.csv};
            \addlegendentry{\scriptsize $v_{i,top}$}
            \addplot[mark=none, dashed, mark size = 1.5pt, line width=0.75pt, color=gray] table [each nth point=1, x={t}, y={vinn}, col sep=comma] {data_simulator_new/ADC_S3_waveforms_veriloga.csv};
            \addlegendentry{\scriptsize $v_{i,bot}$}
            \addplot[mark=none, mark size = 0.5pt, line width=.5pt,  color=gray] table [each nth point=1, x={t}, y={vcm}, col sep=comma] {data_simulator_new/ADC_S3_waveforms_veriloga.csv};
            \addlegendentry{\scriptsize $V_{cm}$}
    
        \end{axis}
        \end{tikzpicture}
        \caption{Behavioral simulation of the proposed SAR ADC with a sinusoidal input signal. The displayed waveforms show the transient response of the CDAC ($v_{top}$, $v_{bot}$), the input signals ($v_{i,top}$, $v_{i,bot}$) and common-mode voltage ($V_{cm}$) for a few conversion cycles. The clock frequency is set at \SI{160}{\MHz}, with $4$ clock cycles allocated to sampling, $2$ for idle states, and $10$ for the actual switching operation. This configuration yields a sampling rate of \SI{10}{\,MSPS}.}
        \label{fig:comp_inputs}
    \end{figure}


    \begin{figure}[t] \centering
       \begin{tikzpicture}
        \begin{axis}[width=16cm, height=10cm, ylabel={\textit{PSD (\SI{}{\dB})}}, xlabel={\textit{Frequency (\SI{}{\MHz})}}, grid=major, enlarge x limits=false, every axis y label/.style={at={(ticklabel cs:0.5)},rotate=90,anchor=near ticklabel}, x filter/.code={\pgfmathdivide{#1}{1000000}}]
            \addplot[mark=none, mark size = 1.5pt, line width=0.7pt, color=gray] table [each nth point=1, x={freq}, y={db20}, col sep=comma] {data_simulator_new/ADC_S3_psd_norm_veriloga.csv};
        \end{axis}
        \end{tikzpicture}
        \caption{Power spectral density (PSD) of the proposed 10-bit SAR ADC with a sinusoidal input signal of \SI{0.395}{\V} and \SI{\sim302.74}{\kHz} frequency. The simulated SQNDR is 62 dB, aligning with an ideal 10-bit resolution. The sampling rate is \SI{10}{\,MSPS} with 1024 FFT points.}
        \label{fig:psd}
    \end{figure}

\end{block}


\end{column}

%%%%%%%%%%%%%%%%%%%%%%%%%%%%%%%%%%%%%%%%%%%%%%%%%%%%%%%%%%%%%%%%%%%%%%%%%%%%%%%%%%%%%%%%%%
\separatorcolumn
%%%%%%%%%%%%%%%%%%%%%%%%%%%%%%%%%%%%%%%%%%%%%%%%%%%%%%%%%%%%%%%%%%%%%%%%%%%%%%%%%%%%%%%%%%

\begin{column}{\colwidth}

    \def\datasetED{datasets_/ED_expanded.csv}
    
    \begin{block}{Comparison}

    The section compares the proposed switching scheme of \Cref{fig:scheme} with state-of-the-art 10-bit SAR ADC solutions. \label{sec:comp}\Cref{tab:comp} showcases the performance metrics against techniques with low energy switching: $E/(C V^2)<32$. The comparison specifically emphasizes architectures utilizing binary weighted configurations without the use of bridge capacitor techniques. Similar to other schemes presented in this analysis, parasitic capacitors and energy consumed during the reset phase are not considered. The conventional SAR is used as the standard reference. \Cref{tab:comp} displays the expected switching energy of multiple switchig techniques and \Cref{fig:ed} shows $E(D)$, i.e., the energy for each digital output code.
    % highlights the comparison of a 10-bit SAR ADC across different techniques and with the actual accumulated switching energy,
    
    \vspace{0.2cm}
    \begin{table}[t] \small
    \caption{Comparison with recent solutions for 10-bit SAR ADC.}
    \begin{threeparttable}
    \centering
    \begin{tabular}{|c||c|c|c|c|c|c|} \hline
                      &\shortstack{Avg. Energy\\$E/(CV^2)$} &\shortstack{$E$ Saving\\(\%)} &\shortstack{Capacitor\\DAC} &\shortstack{Area Red.\\(\%)} &\shortstack{$V_{\text{cm}}$\\Variation} &\shortstack{\#\\Ref.}  \\ \hline
    \cite{Zhu13}      & 31.88     & 97.66       & 512         & 75          & 0.5Vref (High)       & 3\\ \hline
    \cite{Sanyal13}   & 21.30     & 98.40       & 512         & 75          & 0.75Vref (Very high)  & 2\\ \hline
    \cite{Wu17}       & 21.30     & 98.40       & 512         & 75          & 0.5LSB (Low)         & 2\\ \hline
    \cite{Li20}       & 21.08     & 98.45       & 512         & 75          & 0.5LSB (Low)         & 3\\ \hline
    \cite{Hu22}       & 21.08     & 98.45       & 512         & 75          & 0.5LSB (Low)         & 3\\ \hline
    \cite{Akbari19}   & 0         & 100         & 1532       & 25.19          & $< 0.25$Vref (High)        & 3\\ \hline
    \cite{Yousefi18}  & 0\,\tnote{1}         & 100         & 248        & 87.89 & 0               & 3\\ \hline
    \textbf{This} &\textbf{-0.12549} &\textbf{100} &\textbf{512} &\textbf{75} &\textbf{0.5LSB (Very low)}\,\tnote{2} &\textbf{3}\\ \hline
    \end{tabular}
        \begin{tablenotes}
            \item [1] \footnotesize{It requires a more complex digital logic implementation and additional clock cycles during conversion.} \\
            \item [2] \footnotesize{$0.5$ LSB in the $2$ most extreme cases out of $2^{N-1}$ cases, $0$ LSB in the rest. No additional clocks.}
        \end{tablenotes}
    \vspace{0.1cm}
    \label{tab:comp}
    \end{threeparttable}
    \end{table}


     \begin{figure}[t] \centering
       \begin{tikzpicture}
    	   % red, blue, green, darkgray, orange, purple, brown, black, gray, olive, teal
        \begin{axis}[width=16.75cm, height=11cm, legend style={at={(0,1.0)}, legend columns=6, anchor=south west}, ylabel={\textit{Switching Energy ($E / C V^2$)}}, xlabel={\textit{Output Code ($D$)}}, grid=major, enlarge x limits=false, every axis y label/.style={at={(ticklabel cs:0.5)},rotate=90,anchor=near ticklabel}, mark size = 1.5pt ]
    
            \addplot[mark=pentagon, mark size = 1.5pt, color=gray] table [each nth point=7, x=D, y={Chen}, col sep=comma] {\datasetED}; \addlegendentry{\scriptsize\cite{Chen14}}
            %
            \addplot[mark=star, mark size = 1.5pt, color=gray] table [each nth point=7, x=D, y={Zhu}, col sep=comma] {\datasetED}; \addlegendentry{\scriptsize\cite{Zhu10}}
            %
            \addplot[mark=o, mark size = 1.5pt, color=gray] table [each nth point=7, x=D, y={Rahimi}, col sep=comma] {\datasetED}; \addlegendentry{\scriptsize\cite{Rahimi14}}
            %
            \addplot[mark=otimes, mark size = 1.5pt, color=gray] table [each nth point=7, x=D, y={Sanyal}, col sep=comma] {\datasetED}; \addlegendentry{\scriptsize\cite{Sanyal13b}}
            %
            \addplot[mark=x, mark size = 1.5pt, color=gray] table [each nth point=7, x=D, y={Yuan11}, col sep=comma] {\datasetED}; \addlegendentry{\scriptsize\cite{Yuan12}}
            %
            \addplot[mark=oplus, mark size = 1.5pt, color=gray] table [each nth point=7, x=D, y={ZZhu}, col sep=comma] {\datasetED}; \addlegendentry{\scriptsize\cite{Zhu13}}
            %
            \addplot[mark=asterisk, mark size = 1.5pt, color=gray] table [each nth point=7, x=D, y={Li}, col sep=comma] {\datasetED}; \addlegendentry{\scriptsize\cite{Li20}}
            %
            \addplot[mark=diamond, mark size = 1.5pt, color=gray] table [each nth point=7, x=D, y={Sanyal_98}, col sep=comma] {\datasetED}; \addlegendentry{\scriptsize\cite{Sanyal13}}
            %
            \addplot[mark=square, mark size = 1.25pt, color=gray] table [each nth point=20, x=D, y={Akbari}, col sep=comma] {\datasetED}; \addlegendentry{\scriptsize\cite{Akbari19}}
            %
            \addplot[mark=|, mark size = 1.25pt, color=gray] table [each nth point=20, x=D, y={Yousefi}, col sep=comma] {\datasetED}; \addlegendentry{\scriptsize\cite{Yousefi18}}
    
            \addplot[mark=none, mark size = .75pt, color=black, line width=2.5pt] table [each nth point=7, x=D, y={S3}, col sep=comma] {\datasetED}; \addlegendentry{\textbf{\scriptsize{This work}}}
    
        \end{axis}
        \end{tikzpicture}
        \caption{$E(D)$ comparison with state-of-the-art 10-bit SAR ADCs.}
        \label{fig:ed}
    \end{figure}

  \end{block}



  % \begin{alertblock}{alert}
  % \end{alertblock}

  \begin{block}{Summary and future work}
      A simple SAR ADC with zero average switching energy has been introduced. The proposed scheme keeps a constant common-mode voltage, minimizes the capacitor area and requires only one clock per transition.
      Each individual cell of the SAR is prepared for implementation at the transistor level, and as part of future development, exploration of a 28nm CMOS process is underway. 
      It has been demonstrated that by carefully selecting connections at the capacitor terminals a simple architecture achieves a zero power consumption during the successive approximations conversion phase.
  \end{block}

  % \begin{block}{Acknowledgements}
  % \end{block}

  \begin{block}{References}

    % \nocite{*}
    % % \footnotesize{\bibliographystyle{plain}\bibliography{poster}}
    \footnotesize{\printbibliography}

  \end{block}

\end{column}

%%%%%%%%%%%%%%%%%%%%%%%%%%%%%%%%%%%%%%%%%%%%%%%%%%%%%%%%%%%%%%%%%%%%%%%%%%%%%%%%%%%%%%%%%%
\separatorcolumn
%%%%%%%%%%%%%%%%%%%%%%%%%%%%%%%%%%%%%%%%%%%%%%%%%%%%%%%%%%%%%%%%%%%%%%%%%%%%%%%%%%%%%%%%%%

\end{columns}
\end{frame}

\end{document}
